\label{chap:conclusions}

This study successfully introduced and evaluated the \textbf{TransferKAN} architecture for the automated classification of human gunshot wounds. By merging the powerful spatial feature extraction capabilities of a pre-trained ResNet152 backbone with the parameter-efficient, non-linear modeling of a Kolmogorov-Arnold Network (KAN), we established a new benchmark for deep learning applications in forensic pathology.

The analysis, conducted on a realistic and imbalanced dataset of 2,551 crime scene images (FDCPUnBGunshotDB), demonstrated that the TransferKAN model significantly improves classification robustness. Specifically, in distinguishing between Entry and Exit wounds, the proposed architecture achieved a weighted Area Under the Curve (AUC) of 0.908 under a strict case-based 5-fold cross-validation scheme. This performance decisively outperformed the baseline ResNet152 classifier reported in prior literature (AUC of 0.821), proving that KANs are superior in handling minority classes and non-linear feature boundaries. 

For the Medical-Legal Shooting Distance (MLSD) classification, despite the inherent difficulties of extreme class imbalance and the stringent cross-validation protocol preventing data leakage, the TransferKAN model yielded a higher Macro AUC (0.749) and Weighted AUC (0.785) compared to traditional CNN baselines. These results strongly confirm our hypothesis that replacing the standard Multi-Layer Perceptron (MLP) with a KAN enhances the model's capacity to discriminate complex patterns in forensic images.

In conclusion, our approach utilizing TransferKAN with k-fold validation obtained the best performance and overall results, demonstrating its viability as a reliable supportive tool for forensic experts. Future work should focus on expanding the dataset with multi-institutional data, exploring the interpretability of the KAN splines to provide explainable AI insights to forensic pathologists, and investigating the integration of multimodal data (such as autopsy reports) into the classification pipeline.
